\section{Algebraic Structures}

We will soon define numerical structures on which the operations of addition and
multiplication are defined. In general, we want to define what is meant by an
\emph{operation}.
\mymargin{operation}%
\index{operation}%

An \emph{operation} on the set $A$ is a function $f\colon A\to(A\to A)$. Given
$a,b\in A$, we have that $f(a)$ is a function $A\to A$, and thus $f(a)(b)$ will
be an element of $A$. Typically, operations are denoted by symbols such as $+$,
$\cdot$, $/$, $\times$, etc. If $*$ is an operation, instead of writing
$*(a)(b)$, we will write $a*b$.

Note that a function $f\colon A\to (A\to A)$ is equivalent to a function
$g\colon A\times A \to A$ by setting $g(a,b) = f(a)(b)$. Thus, one can
equivalently think of an operation as a function of two variables.

\begin{example}
The usual addition operation $+$ defined on numerical sets is an operation. If
$\NN$ is the set of natural numbers (which we will define shortly), we can think
of $+\colon \NN\to (\NN\to \NN)$ as a function that associates to a number $n\in
\NN$ the function $+n$ that increments its argument by $n$: $(+3)(5)=3+5=8$.
\end{example}

When operation symbols are used in an expression, the problem of establishing
the precedence of operations arises. For example, if we write $a*b*c$, we must
decide whether we mean $(a*b)*c$ (left associativity) or $a*(b*c)$ (right
associativity). If $(a*b)*c=a*(b*c)$, we will say that the operation is
associative, in which case it is obviously unnecessary to insert parentheses. If
parentheses are not inserted, it is usually understood that the association is
to the left. For example, $5-4-3 = (5-4)-3$ and not $5-(4-3)$. In some rare
cases, associativity is to the right; a typical example is exponentiation:
$2^{3^2} = 2^{(3^2)} \neq (2^3)^2$. The reason is that $(a^b)^c$ is more easily
written as $a^{b\cdot c}$. In programming languages that implement the
exponentiation operation, one must be very careful: in some languages (e.g.,
Fortran), exponentiation associates to the left like other operations, while in
other languages (e.g., Python), it associates to the right.

When different operations appear in the same expression, some operations take
precedence over others. For example, in $a+b\cdot c$, multiplication takes
precedence over addition, so the expression should be interpreted as $a+(b\cdot
c)$, not as $(a+b)\cdot c$. Addition and subtraction have the same precedence
and should therefore be interpreted from left to right: $a-b+c = (a-b)+c$. For
division, we will rarely use the symbol $/$; in that case, division associates
to the left and has the same precedence as multiplication: $a/b/c=(a/b)/c$,
$a/b\cdot c=(a/b)\cdot c$. This rule is sometimes counterintuitive, so we will
avoid writing expressions of this type, especially because the multiplication
dot can be omitted, as in $a/bc$, making the expression even more ambiguous.
When in doubt: use parentheses! To denote division, we will more often use the
fraction symbol $\frac a b$, which does not require specifying associativity and
precedence since the length of the fraction line determines how the operands are
associated:
\[
  \frac{\,\frac a b\,}{c}=(a/b)/c, \qquad 
  \frac{a}{\,\frac bc\,}=a/(b/c).
\]  

The numerical sets we aim to define (natural, integer, rational, real, complex
numbers) are all equipped with the two operations of addition and
multiplication, as well as (except for complex numbers) an order relation. Based
on the properties that these operations satisfy, they can fall into the
following abstract algebraic structures.

For now, we will list the structures that we will encounter later. As we
identify examples of sets that fall into these structures, the meaning of these
definitions will become clearer.

\begin{definition}[Monoid]%
  \label{def:monoide}%
    Let $A$ be a set on which an operation $*$ is defined. If the operation is
  associative: $x*(y*z) = (x*y)*z$ and has an identity element $e\in A$ such
  that $e*x = x*e = x$, we will say that $A$ is a \emph{monoid}. If, in
  addition, the operation is commutative, $x*y=y*x$, we will say that $A$ is a
  \emph{commutative monoid} (or abelian).
  
    When the symbol of addition $+$ is used for the operation $*$, we will say
  that the monoid is additive. In this case, the identity element is called
  \emph{zero} and is denoted by $0$. If instead we use the symbol of
  multiplication $x\cdot y$, we will say that the monoid is multiplicative. In
  this case, the identity element is called \emph{unity} and is usually denoted
  by $1$.
\end{definition}
  
\begin{example}
The set of natural numbers is a commutative monoid both with the addition
operation $+$ (additive monoid) and with the multiplication operation $\cdot$
(multiplicative monoid).
\end{example}
  
\begin{definition}[Group]
  \label{def:gruppo}%
    A set $G$ on which an \emph{operation} $*$ is defined is said to be a
  \emph{group}%
\mymargin{group}%
\index{group} if the operation has the following properties:
  \begin{enumerate}
    \item Associative: $\forall x,y,z\in G\colon (x*y)*z = x*(y*z)$;
    \index{property!associative}%
    \index{associativity}%
    \item Existence of identity element: 
    \index{element!identity}%
    \index{identity}%
    $\exists e\in G\colon \forall x\in G \colon e*x=x*e = x$;
    \item Existence of inverse: 
    \index{element!inverse}%
    \index{inverse}%
    $\forall x\in G\colon \exists y\in G\colon x*y=y*x=e$.
  \end{enumerate}
  Additionally, the group is said to be \emph{abelian}%
\mymargin{abelian group}%
\index{abelian} or \emph{commutative} if the property holds:
  \begin{enumerate}
    \item[4.] Commutative: $\forall x,y\in G\colon x*y = y*x$.
    \index{property!commutative}%
    \index{commutativity}% 
  \end{enumerate}
  
    When the operation is denoted by the symbol $+$ (addition), we will say that
  the group is additive, and we will denote the identity element by $0$
  \index{zero}%
    (zero), and the inverse of $x$ will be called \emph{opposite} and denoted by
  $-x$.
  \index{opposite}%
    If instead the symbol $\cdot$ (multiplication) is used, we will say that the
  group is multiplicative, the identity element may be denoted by the symbol $1$
  (one or unity), and
  \index{one}\index{unity}%
    the inverse of $x$ may be called \emph{reciprocal} and is usually denoted by
  $x^{-1}$ or $\frac 1 x$.
  \index{reciprocal}%
\end{definition}
  
The identity element of a group is unique. If $x$ and $y$ were two identity
elements, we would have $x = x*y = y$. The inverse is also unique: indeed, if
$y$ and $z$ were two inverses of $x$, we would have $y = y * x * z = z$.

\begin{example}
    The set of integers $\ZZ$ (which we will define later) with the addition
  operation is an abelian group. The identity element is zero, and the opposite
  of $x$ is $-x$. On $\ZZ$, we will also define the multiplication operation,
  but $\ZZ$ with the multiplication operation is not a group because not all
  elements have a multiplicative inverse in $\ZZ$.
\end{example}

\begin{example}[Symmetric Group]
    If $X$ is any set, one can consider the set $X!$ of bijective functions on
  $X$:
  \[
    X! = \ENCLOSE{f\colon X\to X\colon f \text{ bijective}}.
  \]
    When $X$ is a finite set, the bijective functions on $X$ are also called
  \emph{permutations} because they permute the elements of $X$.

    As an exercise, one can verify that the operation $\circ$ of composition
  makes $X!$ a group. This group is called the \emph{symmetric group} of $X$ (it
  is also usually denoted by $S_X$). The symmetric group is not abelian if $X$
  has more than two elements.
\end{example}

\begin{definition}[Ring and Field]
  \label{def:anello}%
  \label{def:campo}%
    Let $A$ be a set on which two operations are defined: addition and
  multiplication. We will say that $A$ is a \emph{ring} if $A$ is an abelian
  group with respect to addition, if multiplication is associative $(x\cdot
  y)\cdot z = x\cdot (y\cdot z)$, and if the distributive property holds
  $x\cdot(y+z) = x\cdot y + x\cdot z$, $(x+y)\cdot z = x\cdot z + y\cdot z$.

    If, in addition, there exists an element $1\in A$ that is the identity for
  multiplication, we will say that $A$ is a ring \emph{with unity}.

    If multiplication is commutative, we will say that $A$ is a
  \emph{commutative ring}.

    If $A$ is a ring and $A\setminus \ENCLOSE{0}$ turns out to be an abelian
  group for the multiplication operation (thus there is a unity $1\neq 0$ and
  every non-zero element has a multiplicative inverse), we will say that $A$ is
  a \emph{field}.
\end{definition}

\begin{example}
    The set of integers $\ZZ$ (which we will define shortly) is a commutative
  ring with unity but is not a field. The set of rational numbers $\QQ$ (which
  we will define shortly) is a field.
\end{example}

A ring $A$ is an additive group, so in rings, the identity element of addition
is denoted by $0$. The additive inverse of $a\in A$ is called the
\emph{opposite} and is denoted by $-a$. Thus, the operation of subtraction is
also defined: $a-b = a+(-b)$.

The same holds in a field $R$ where, in addition, every non-zero element also
has a multiplicative inverse. If $x\in R\setminus\ENCLOSE{0}$, its
multiplicative inverse is called the \emph{reciprocal} and is denoted by $1/x$
or $x^{-1}$. Thus, the operation of division is also defined: $a/b = a\cdot
(1/b)$ if $b\neq 0$.

In general, if $A$ is a ring, the familiar properties hold:
\[
  0\cdot x = x\cdot 0 = 0, \qquad
  (-1)\cdot x = x \cdot (-1) = -x.
\]
To prove the first, we have:
\[
  0\cdot x = 0\cdot x + x + (-x) = (0+1)\cdot x + (-x) = x + (-x) = 0
\]
and writing the addends in the opposite order, we also obtain $x\cdot 0 = 0$. We
will say that $0$ is an \emph{absorbing element}
\mymargin{absorbing element}%
\index{element!absorbing}%  
\index{absorbing}%
for multiplication. Then we can also prove the second property:
\[
   (-1)\cdot x = (-1)\cdot x + x + (-x) = (-1 + 1)\cdot x + (-x) = 0 + (-x) = -x
\]
and again, writing the addends in the opposite order, we obtain $x\cdot(-1)=-x$.

If $A$ is a field, the \emph{zero-product property} holds:
\index{zero-product}%
\mymargin{zero-product}%
\[
  x\cdot y = 0 \iff x=0 \lor y=0.
\]
The implication to the left is the absorbing property, which we have already
verified. Conversely, if $x\cdot y = 0$ and $x\neq 0$, then we can multiply both
sides by $x^{-1}$ to obtain $y=x^{-1}\cdot 0 = 0$.

\begin{definition}[Totally Ordered Group]
  \label{def:gruppo_ordinato}%
  \label{def:campo_ordinato}%
  We will say that $G$ is a \emph{totally ordered group}
\mymargin{totally ordered group}%
\index{group!ordered}% 
    if $G$ is a group (with operation $*$), if it is also a totally ordered set
  (with relation $\le$), and if the group operation preserves the ordering,
  i.e., the property holds:
  \begin{enumerate}
    \item[1.] Monotonicity:
    if $x\le y$, then for every $z$, we have:
      \[
      x*z \le y*z \qquad\text{and}\qquad z*x \le z*y.
      \] 
  \end{enumerate}
    We will say that the totally ordered group is \emph{dense} or
  \emph{continuous} if the order relation also has such properties.

  We will say that $R$ is an \emph{ordered field}
\mymargin{ordered field}%
\index{field!ordered}%
    if it is a field, if with respect to addition it is a totally ordered
  abelian group (thus addition preserves the ordering $x\le y \implies x+z\le
  y+z$), and if, in addition, multiplication by a positive number preserves the
  ordering:
  \begin{enumerate}
    \item[2.] Monotonicity: 
    if $x\le y$ and $0\le z$, then $x\cdot z \le y\cdot z$. 
  \end{enumerate} 

    We will say that the ordered field is \emph{continuous} if the ordering is
  continuous.
\end{definition}

\begin{example}
We will see that $\ZZ$ (the set of integers) is an example of a totally ordered
additive group but not dense.
\end{example}

\begin{example}
We will see that $\QQ$ (the set of rational numbers) is an example of a dense
ordered field but not continuous, while $\RR$ (the set of real numbers) is (in a
certain sense the only) example of a continuous ordered field.
\end{example}

\begin{example}
A trivial example of a dense and continuous totally ordered group is
$X=\ENCLOSE{0}$.
\end{example}

If $G$ is a totally ordered additive group with identity element $0$, we have:
\[
  x\le y  \iff 0 \le y-x.
\]
This means that the ordering is uniquely determined by comparison with the
identity element, i.e., by knowing which elements are positive.

If $G$ is a totally ordered additive group, we can express a monotonicity
property for strict inequalities: if $x<y$, then $x+z<y+z$ and $z+x<z+y$. This
is to observe that if the equality $x+z=y+z$ holds, by adding $-z$ to both
sides, we obtain $x=y$.

If $R$ is an ordered field and we take $x\in R$, we can observe that $x\ge 0$ if
and only if $-x\le 0$. Indeed, it suffices to add $-x$ to the first inequality
to obtain the second and vice versa to add $x$ to the second to obtain the first
(monotonicity of addition).

We now verify the validity of the sign rule for the product: if $x\ge 0$ and
$y\ge 0$, then $x\cdot y \ge 0$, $(-x)\cdot (-y) \ge 0$, $(-x)\cdot y = x \cdot
(-y) \le 0$. The first property, $x\cdot y\ge 0$, is given by hypothesis
(monotonicity property of the product applied to $y\ge 0$ by multiplying both
sides by $x\ge 0$). Consequently, $-(x\cdot y)\le 0$, i.e., $(-x)\cdot y \le 0$
and $x\cdot (-y) \le 0$. Changing the sign again on this last inequality, we
obtain $(-x)\cdot(-y)\ge 0$.

Thus, we can deduce that for every $x\in R$, we have $x^2=x\cdot x\ge 0$ because
whether $x\ge 0$ or $x\le 0$, the product will always be $\ge 0$.

In particular, we can deduce that $1=1\cdot 1\ge 0$. And since by hypothesis of
the field, $1\neq 0$, we must have $1>0$.

\begin{exercise}
Let $R$ be an ordered field, $m\in R$, $m>0$. Then $m^{-1}>0$. Moreover (if
$m>0$) for every $x,y\in R$, we have:
\[
  x<y \iff m\cdot x < m\cdot y.
\]
\end{exercise}

\begin{exercise}
Prove that every ordered field is dense.
\end{exercise}
\begin{proof}[Solution.]
Given $x,y\in R$ with $x<y$, we must find $z\in R$ such that $x<z<y$. If $x<y$,
we have $x+x<x+y<y+y$, which we can rewrite as
\[
(1+1)\cdot x < x+y < (1+1)\cdot y.
\]
Setting $2=1+1$, multiplying by the inverse of $2$ (which is positive), we
obtain
\[
x < \frac{x+y}{2} < y.
\]
\end{proof}